\documentclass[11pt,a4paper]{article}
\usepackage[T1]{fontenc}
\usepackage[utf8]{inputenc}
\usepackage{lmodern}
\usepackage{amsmath,amssymb,amsthm}
\usepackage[margin=1.15in]{geometry}
\usepackage[hidelinks]{hyperref}
\emergencystretch=3em

\newtheorem{theorem}{Theorem}[section]
\newtheorem{proposition}[theorem]{Proposition}
\newtheorem{lemma}[theorem]{Lemma}
\newtheorem{corollary}[theorem]{Corollary}
\theoremstyle{definition}
\newtheorem{definition}[theorem]{Definition}
\newtheorem{remark}[theorem]{Remark}

\newcommand{\A}{\mathcal{A}}
\newcommand{\K}{\mathcal{K}}
\newcommand{\C}{\mathcal{C}}
\newcommand{\V}{\mathcal{V}}
\newcommand{\Po}{\mathcal{P}}
\newcommand{\Z}{\mathbb{Z}}
\newcommand{\Q}{\mathbb{Q}}
\newcommand{\R}{\mathbb{R}}
\newcommand{\F}{\mathbb{F}}
\newcommand{\colim}{\operatorname{colim}}
\newcommand{\aug}{\operatorname{aug}}
\newcommand{\Aut}{\operatorname{Aut}}
\newcommand{\Hom}{\operatorname{Hom}}
\newcommand{\id}{\mathrm{id}}

\title{Exactness of colimits and amenability are independent for categories
of finite substructures\thanks{AI-produced. The results are due to a run of
GPT-6 Astra (OpenAI, ChatGPT Pro mode); the text was written, and every proof
derived again, by Claude Fable~5.2 (Anthropic) at effort max; see
Section~\ref{sec:provenance} for who proved, wrote and checked what. No human
has checked this text. To the extent that any rights exist in it, they are
dedicated to the public domain under CC0~1.0.}}
\author{}
\date{25 November 2026}

\begin{document}
\maketitle

\begin{abstract}
Let $M$ be a countable homogeneous structure, $G$ its automorphism group and
$\A$ the category of its finite substructures and embeddings. The colimit
functor on diagrams of $R$-modules indexed by $\A$ is exact if and only if
the functor of coinvariants $N\mapsto N_G$ is exact on the $R[G]$-modules in
which every element has an open stabilizer; this is a property of the topos
of continuous $G$-sets. We compare it with amenability of $G$, the existence
of invariant probability measures on compact spaces on which $G$ acts, and
show that neither implies the other. For the generic partially ordered set
with an independent linear order, colimits are exact over every ring, while
the group is not amenable by a result of Kechris and Soki\'c. For the vector
space of countable dimension over the field with two elements, with a generic
linear order on its non-zero vectors, the group is amenable and colimits are
exact over the rational numbers, but they are exact over no ring in which
$2$ is not invertible. The first statement rests on a transfer theorem: if a
subgroup $H$ has exact coinvariants and finitely many orbits on each
$G/G_B$, then exactness for $G$ over a field $k$ is equivalent to the
existence of a $G$-invariant, finitely additive, $k$-valued function of total
mass $1$ on the expansions that $H$ defines; positivity is not required. For
posets this function has integer values, given by Euler characteristics of
cells of order cones, which are values of the measure of Harman and Snowden
for the automorphism group of the ordered real line. A consequence is that
the separation property proved by Evans for Ramsey structures holds for a
structure that is not a Ramsey structure but is obtained from one by
forgetting a relation.
\end{abstract}

\section{Introduction and results}\label{sec:intro}

\paragraph{The setting.} We use the setting of \cite[\S3.1]{D1}. Let $M$ be a non-empty countable set and $G$ a group of permutations of $M$. For $B\subseteq M$ let
$G_B$ be the pointwise stabilizer of $B$. A subgroup of $G$ is \emph{open} if
it contains $G_B$ for some finite $B$; this is the topology of pointwise
convergence on $G$. Let $\K$ be a set of finite subsets of $M$ such that
\begin{enumerate}
\item[\textup{(K1)}] $\varnothing\in\K$; $gB\in\K$ for $B\in\K$ and $g\in G$;
every finite subset of $M$ is contained in a member of $\K$.
\end{enumerate}
The category $\A=\A(M,G,\K)$ has the non-empty members of $\K$ as objects,
and its arrows $B\to C$ are the maps $g|_B$ with $g\in G$ and $gB\subseteq C$.
If $M$ is the Fra\"iss\'e limit of a class of finite structures, $G=\Aut(M)$
and $\K$ is the set of all finite subsets, then every embedding between
finite substructures of $M$ extends to an automorphism, so $\A$ is the
category of non-empty finite substructures of $M$ and embeddings, and it is
equivalent to the category of all non-empty structures of the class with
embeddings. In Section~\ref{sec:vector} the members of $\K$ are the finite
subspaces of a vector space, without the zero vector.

All rings are commutative with $1$. For a ring $R$, an $R[G]$-module $N$ is
\emph{smooth} if every element of $N$ is fixed by an open subgroup. Its
\emph{coinvariants} are $N_G=N/I_GN$, where $I_GN$ is the submodule generated
by the elements $gn-n$; the functor $N\mapsto N_G$ is right exact.

\begin{definition}\label{def:exact}
The category $\A$ is \emph{exact over $R$} if the functor
$\colim_{\A}$ from functors $\A\to R\text{-Mod}$ to $R$-modules is exact.
\end{definition}

Colimits of modules are computed on the underlying abelian groups, so a
category that is exact over $\Z$ is exact over every ring; and exactness is
invariant under equivalence of categories.

The group $G$ is \emph{amenable} if every continuous action of $G$ on a
non-empty compact Hausdorff space has an invariant regular Borel probability measure, and \emph{extremely amenable} if every such action has a fixed
point.

\begin{theorem}\label{thm:dict}
Assume \textup{(K1)}. The category $\A(M,G,\K)$ is exact over $R$ if and only
if the functor of coinvariants preserves monomorphisms of smooth
$R[G]$-modules. Equivalently: in the topos of continuous $G$-sets, the left
adjoint of the functor that gives an $R$-module the trivial action of $G$ is
exact on $R$-module objects. In particular exactness over $R$ depends only
on this topos.
\end{theorem}

The first sentence is due to the run described in
Section~\ref{sec:provenance}; the functors in its proof are those of Artin's
theorem as written out for modules by Di, Li, Liang and Xu
\cite[Theorem~5.1]{DLLX}. Section~\ref{sec:dict} also proves a criterion in
terms of the category alone: $\A$ is exact over $R$ if and only if every
finite family of parallel arrows is equalized by a combination of arrows
with coefficients in $R$ of sum $1$ (Proposition~\ref{prop:affine}); this is
the criterion of Isbell and Mitchell \cite{IM} in the present setting.

\begin{theorem}\label{thm:transfer}
Assume \textup{(K1)}. Let $H$ be a subgroup of $G$ with finitely many orbits
on $G/G_B$ for every $B\in\K$, and let $k$ be a ring such that coinvariants
preserve monomorphisms of smooth $k[H]$-modules (smooth for $H$: every element is fixed by $H\cap G_B$ for some finite $B$). For $B\in\K$ let
$E(B)=G_B\backslash G/H$.
\begin{itemize}
\item[\textup{(a)}] If there are functions $\mu_B\colon E(B)\to k$, $B\in\K$,
of total sum $1$, compatible with the restriction maps $E(C)\to E(B)$ for
$B\subseteq C$ and with the action of $G$ \textup{(}a \emph{normalized
invariant measure}, Definition~\textup{\ref{def:measure})}, then
$\A(M,G,\K)$ is exact over every $k$-algebra.
\item[\textup{(b)}] If $k$ is a field and $\A(M,G,\K)$ is exact over $k$,
there is such a family.
\item[\textup{(c)}] If $G$ is amenable and $k\supseteq\Q$, there is such a
family; so $\A(M,G,\K)$ is then exact over every $\Q$-algebra.
\end{itemize}
\end{theorem}

When $H$ is the automorphism group of an expansion $M^*$ of $M$, the set
$E(B)$ is the set of expansions of $B$ that occur in $M^*$, and a normalized
invariant measure with values in $[0,1]$ is what is called a consistent
random expansion; for excellent pairs of classes its existence is the
starting point of Zucker's criterion for amenability
\cite[\S3 and Theorem~3.1]{Zu}. Theorem
\ref{thm:transfer} says that exactness is the same linear condition without
positivity (Remark~\ref{rem:zucker}).

\begin{theorem}\label{thm:posets}
Let $\Po$ be the class of finite structures $(P,\prec,<)$, where $\prec$ is a
strict partial order and $<$ a linear order on $P$, let $M_{\Po}$ be its
Fra\"iss\'e limit and $G=\Aut(M_{\Po})$. The category of non-empty members of
$\Po$ and embeddings is exact over every ring. The group $G$ is not amenable
\textup{(}Kechris and Soki\'c \cite[Remark~3.2]{KS}\textup{)}.
\end{theorem}

\begin{theorem}\label{thm:vector}
Let $\V$ be the category of finite-dimensional non-zero vector spaces over
$\F_2$ with a linear order on the set of their non-zero vectors, and of
injective linear maps that preserve the order. Let $(V_\infty,<)$ be its
Fra\"iss\'e limit and $G$ the group of linear automorphisms of $V_\infty$ that
preserve the order.
\begin{itemize}
\item[\textup{(a)}] $G$ is amenable.
\item[\textup{(b)}] $\V$ is exact over every $\Q$-algebra.
\item[\textup{(c)}] $\V$ is not exact over any ring in which $2$ is not
invertible; in particular not over $\F_2$ and not over $\Z$.
\end{itemize}
\end{theorem}

So exactness over $\Z$ does not imply amenability, and amenability does not
imply exactness over $\Z$ or over a field of characteristic $2$. Extreme
amenability does imply exactness over $\Z$: this is Lemma~3.7\,(i) of Evans
\cite{Ev}, in dual form (Remark~\ref{rem:evans}).

The second statement alone has an elementary witness with non-rigid objects.
For the category of finite non-empty sets and injections the group is the
symmetric group of a countably infinite set. It is amenable by
Lemma~\ref{lem:Re}: the uniform distribution $\nu$ on the bijections of $Q$
satisfies $\nu a=\nu a'$ for all arrows $a,a'\colon P\to Q$. The average of
the bijections of $Q$ is an affine equalizer over $\Q$ of all arrows
$P\to Q$, and an affine equalizer of the automorphisms of an $n$-element set
has augmentation in $n!R$; so this category is exact over $R$ exactly when
$R$ is a $\Q$-algebra (Proposition~\ref{prop:affine}). In
Theorem~\ref{thm:vector} every object is rigid and the category is exact
over $\Q$. We do not know whether $\V$ is exact over a ring in which $2$ is
invertible and which is not a $\Q$-algebra. Amenability does not imply
exactness over $\Q$ in general: for the group of translations of $M=\Z$,
with $\K$ the set of finite subsets, the group is discrete and amenable,
the objects are rigid, and the inclusion of $\{0\}$ in $\{0,1\}$ and the
arrow $0\mapsto1$ have no affine equalizer over any non-zero ring, because
the coefficients of an equalizer would form a finitely supported function
on $\Z$ that is invariant under translation. By
Theorem~\ref{thm:transfer}\,(c) amenability does imply exactness over $\Q$
when $G$ has a subgroup $H$ as in that theorem.

\begin{corollary}\label{cor:evans}
Let $M=M_{\Po}$, $G=\Aut(M)$, let $W$ be a set with an action of $G$ in which
every point has an open stabilizer, $R$ a ring, $Y$ an $R[G]$-submodule of
the free module $R[W]$, and $x\in R[W]$. Let $U$ be an open subgroup of $G$
fixing $x$, and $\Omega_U\colon R[W]\to R[U\backslash W]$ the map that sends
a point to its $U$-orbit. Then $x\in Y$ if and only if
$\Omega_U(x)\in\Omega_U(Y)$. If $R$ is a field and $x\notin Y$, there is a
$U$-invariant function $f\colon W\to R$ whose linear extension vanishes on
$Y$ and not on $x$.
\end{corollary}

For $\omega$-categorical Ramsey structures, $W$ a sort of $M^{\mathrm{eq}}$ and $U$ the pointwise stabilizer of the support of $x$, these are Theorem~1.4 and Corollary~1.5 of \cite{Ev}. In Remarks~2.12 of \cite{Ev} Evans asks whether
the separation property holds for reducts of Ramsey structures. The
structure $M_{\Po}$ is a reduct of a Ramsey structure and is not one itself
(its group is not amenable); Corollary~\ref{cor:evans} gives the property for
it over every field, and Theorem~\ref{thm:transfer}\,(c) gives it in
characteristic $0$ for every reduct with an amenable automorphism group
(Remark~\ref{rem:reducts}).

\paragraph{What is known and what is new.} Known: the non-amenability in
Theorem~\ref{thm:posets} \cite{KS}; the implication from extreme amenability
to exactness \cite{Ev}; Artin's dictionary between continuous $G$-sets and
sheaves on orbit categories, in the form of \cite{DLLX} and, for Fra\"iss\'e
classes, of Caramello \cite[Theorem~2.9, Proposition~2.10]{Ca}; the criterion
of Isbell and Mitchell \cite{IM}; the Euler-characteristic measure of Harman
and Snowden \cite{HS}, of which the weights of Lemma~\ref{lem:weights} are
values (Remark~\ref{rem:HS}). New, to our knowledge after the search
described in Section~\ref{sec:provenance}: Theorems~\ref{thm:transfer},
\ref{thm:vector}, the exactness in Theorem~\ref{thm:posets} and
Corollary~\ref{cor:evans}. The theory of \cite{HS} does not give these
statements: its exactness results concern modules over a completed group
algebra, and smoothness of a $k[G]$-module in the sense above is weaker than smoothness over that algebra \cite[Remark~10.10]{HS}; and by Snowden \cite[Proposition~8.2]{Sn} the
class of finite posets has no regular measure in their sense.

\section{Coinvariants, the topos and affine equalizers}\label{sec:dict}

Throughout, $(M,G,\K)$ satisfies (K1) and $\A=\A(M,G,\K)$. By (K1) the
non-empty members of $\K$ form a directed set under inclusion, inclusions
are arrows of $\A$, and every arrow $f\colon B\to C$ factors as
$B\xrightarrow{\,g|_B\,}gB\subseteq C$ for any $g\in G$ with $g|_B=f$.

\begin{proof}[Proof of Theorem~\ref{thm:dict}]
For $V\colon\A\to R\text{-Mod}$ let $S(V)=\colim_{B}V(B)$, the colimit over
the directed set of non-empty members of $\K$ and inclusions. Write $[B,v]$
for the class of $v\in V(B)$. The rule $g[B,v]=[gB,V(g|_B)v]$ is well defined,
because $g|_C$ restricts to $g|_B$ for $B\subseteq C$, and it makes $S(V)$ an
$R[G]$-module in which $[B,v]$ is fixed by $G_B$; so $S(V)$ is smooth. The
functor $S$ is exact, because directed colimits of modules are exact.

\emph{Claim: $\colim_{\A}V\cong S(V)_G$, naturally in $V$.} The module
$\colim_{\A}V$ is the quotient of $\bigoplus_BV(B)$ by the relations
$v\sim V(f)v$ for arrows $f$. By the factorization above these relations are
generated by those for inclusions, whose quotient is $S(V)$, and those for
the bijections $g|_B\colon B\to gB$, which identify $[B,v]$ with $g[B,v]$.

For a smooth $R[G]$-module $N$ define $T(N)\colon\A\to R\text{-Mod}$ by
$T(N)(B)=N^{G_B}$ and $T(N)(g|_B)n=gn$. This is well defined: another
element of $G$ with the same restriction to $B$ is $gu$ with $u\in G_B$; and
$gn$ is fixed by $G_C$ when $gB\subseteq C$, since
$g^{-1}G_Cg=G_{g^{-1}C}\subseteq G_B$. Since $N$ is smooth,
$S(T(N))=\bigcup_BN^{G_B}=N$ as $R[G]$-modules, and a monomorphism
$N'\to N$ gives a monomorphism $T(N')\to T(N)$.

If coinvariants preserve monomorphisms of smooth modules, then
$\colim_{\A}=(-)_G\circ S$ preserves monomorphisms, and it is right exact;
so $\A$ is exact over $R$. Conversely, if $\A$ is exact over $R$ and
$N'\subseteq N$ are smooth, then
$N'_G=\colim_{\A}T(N')\to\colim_{\A}T(N)=N_G$ is injective.

A continuous $G$-set is a set with an action of $G$ in which every point
has an open stabilizer, so an $R$-module object of the topos of continuous
$G$-sets is a smooth $R[G]$-module. The functor $\gamma^*$ from sets to this
topos gives a set the trivial action, and for a smooth $N$ and an
$R$-module $B$ we have $\Hom_{R[G]}(N,\gamma^*B)=\Hom_R(N_G,B)$. So the left
adjoint of $\gamma^*$ on $R$-module objects is $N\mapsto N_G$. An equivalence
of toposes commutes with $\gamma^*$ up to isomorphism, which is unique, and
induces an equivalence of the categories of $R$-module objects; hence the
last assertion.
\end{proof}

\begin{remark}\label{rem:topos}
Under the equivalence of Artin between continuous $G$-sets and sheaves for
the atomic topology on an orbit category of $G$, the functor $T$ is the
inclusion of sheaves into presheaves and $S$ is the associated sheaf (when the set of fixed points of $G_B$ in $M$ is $B$ for all $B\in\K$, as in the examples of this paper, so that $\A$ is the opposite of the orbit category on the $G/G_B$; in general $S$ is exact and left adjoint to the fully faithful $T$) \cite[Theorem~5.1 and the equivalence displayed after it]{DLLX}. When $M$ is
the Fra\"iss\'e limit of the finitely presentable models of a theory of
presheaf type, the topos of continuous $\Aut(M)$-sets is the classifying
topos of the theory of homogeneous models
\cite[Theorem~2.9, Proposition~2.10]{Ca}. By Theorem~\ref{thm:dict},
exactness of colimits over the category of finite models is then an
invariant of that theory up to Morita equivalence.
\end{remark}

\begin{lemma}\label{lem:open}
If coinvariants preserve monomorphisms of smooth $R[G]$-modules, the same
holds for every open subgroup $U$ of $G$. Then for every smooth
$R[U]$-module $N$ the map $N^U\to N_U$ is injective.
\end{lemma}

\begin{proof}
For a smooth $R[U]$-module $N$ the induced module $R[G]\otimes_{R[U]}N$ is
smooth, because $g\otimes n$ is fixed by $gU_ng^{-1}$ and the stabilizer
$U_n$ of $n$ is open in $G$. Induction is exact, since $R[G]$ is free as a
right $R[U]$-module, and
$(R[G]\otimes_{R[U]}N)_G=R\otimes_{R[U]}N=N_U$. So a monomorphism
$N'\to N$ of smooth $R[U]$-modules gives a monomorphism $N'_U\to N_U$. For
$n\in N^U$ the submodule $Rn$ has trivial action, so $Rn=(Rn)_U\to N_U$ is
injective.
\end{proof}

For objects $P,C$ of $\A$ write $\A(P,C)$ for the set of arrows, which is
finite, and $R[\A(P,C)]$ for the free module on it. An arrow
$a\colon P\to Q$ induces $t\mapsto ta$ from $R[\A(Q,C)]$ to $R[\A(P,C)]$,
and $\aug$ denotes the sum of the coefficients.

\begin{definition}\label{def:affine}
An \emph{affine equalizer over $R$} of arrows $a_0,\dots,a_k\colon P\to Q$ is
an element $t\in R[\A(Q,C)]$, for some object $C$, with $\aug(t)=1$ and
$ta_i=ta_0$ for all $i$.
\end{definition}

\begin{proposition}\label{prop:affine}
The category $\A$ is exact over $R$ if and only if every finite family of
arrows $a_0,\dots,a_k\colon P\to Q$ of $\A$ has an affine equalizer over $R$.
\end{proposition}

For arbitrary small categories with common successors this is the criterion
of Isbell and Mitchell \cite[Theorem~3.1, Proposition~3.3 and the remark after it, for a commutative ring in place of $\Z$]{IM}. We give a
direct proof in the present setting.

\begin{proof}
For an object $P$ let $E_P$ be the free $R$-module on the set of maps
$g|_P$, $g\in G$, with $G$ acting by composition; it is the permutation
module on $G/G_P$, it is smooth, and $(E_P)_G=R$ by the augmentation.

Suppose $\A$ is exact over $R$ and let $a_0,\dots,a_k\colon P\to Q$. Let
$\phi\colon E_Q\to E_P^k$ have the components $t\mapsto ta_i-ta_0$; it is
$G$-linear. Put $I=\phi(E_Q)$. Coinvariants are right exact, so
$I_G=R/\aug(\ker\phi)$. The composite $E_Q\to I\to E_P^k\to(E_P^k)_G=R^k$
is zero, because $\aug(ta_i)=\aug(ta_0)$; so $I_G\to(E_P^k)_G$ is zero. By
Theorem~\ref{thm:dict} it is injective. Hence $I_G=0$, and there is
$t\in\ker\phi$ with $\aug(t)=1$. Its support is a finite set of maps
$Q\to M$, whose images lie in some $C\in\K$; so $t$ is an affine equalizer.

Conversely suppose every finite family has an affine equalizer. Let
$N'\subseteq N$ be smooth and $n'\in N'\cap I_GN$, say
$n'=\sum_{i=1}^k(g_i-1)n_i$ with $g_i\in G$, $n_i\in N$. Choose $P\in\K$
such that $G_P$ fixes $n_1,\dots,n_k$, and $Q\in\K$ containing $P$ and all
$g_iP$. Let $a_0\colon P\to Q$ be the inclusion and $a_i=g_i|_P$, and let
$t=\sum_ht_h[h]\in R[\A(Q,C)]$ be an affine equalizer. For each $h$ choose
$\hat h\in G$ with $\hat h|_Q=h$. For $n\in N^{G_P}$ the rule
$\Phi_n([g|_P])=gn$ defines an $R$-linear map $\Phi_n\colon R[\A(P,C)]\to N$,
and $\Phi_n([ha_i])=\hat hg_in$, because $\hat hg_i$ restricts to $ha_i$ on
$P$. Therefore
\[
 \sum_ht_h\hat hg_in_i=\Phi_{n_i}(ta_i)=\Phi_{n_i}(ta_0)=\sum_ht_h\hat hn_i ,
\]
and summing over $i$ gives $\sum_ht_h\hat hn'=0$. As $\sum_ht_h=1$,
\[
 n'=\sum_ht_hn'-\sum_ht_h\hat hn'=-\sum_ht_h(\hat h-1)n'\in I_GN'.
\]
So $N'_G\to N_G$ is injective, and $\A$ is exact over $R$ by
Theorem~\ref{thm:dict}.
\end{proof}

\section{Ramsey categories, amenability and the transfer}\label{sec:transfer}

\subsection{Ramsey categories}

\begin{definition}\label{def:RP}
A category $\C$ with finite sets of arrows has the \emph{Ramsey property}
if for all objects $P,Q$ and every $r\ge1$ there is an object $C$ such that
for every map $\chi\colon\C(P,C)\to\{1,\dots,r\}$ there is $g\in\C(Q,C)$
with $\chi$ constant on $\{ga:a\in\C(P,Q)\}$. It has property
\textup{(Re)} if for all arrows $a_0,\dots,a_k\colon P\to Q$ and every
$\varepsilon>0$ there are an object $C$ and a probability vector $\nu$ on
$\C(Q,C)$ with $\sum_{i=1}^k\|\nu a_i-\nu a_0\|_1<\varepsilon$.
\end{definition}

Definition~\ref{def:affine} makes sense in any such category.

\begin{lemma}\label{lem:ramsey}
In a category with finite sets of arrows and the Ramsey property, every
finite family $a_0,\dots,a_k\colon P\to Q$ has an affine equalizer over $\Z$,
and property \textup{(Re)} holds.
\end{lemma}

\begin{proof}
\emph{Affine equalizers.} Let $N=3k+1$ and let $C$ be as in
Definition~\ref{def:RP} for $P$, $Q$ and $r=N^k$. Let
$F=\Z[\C(P,C)]^k\oplus\Z$ and let $L\subseteq F$ be the image of
$t\mapsto((ta_i-ta_0)_i,\aug t)$ on $\Z[\C(Q,C)]$. An affine equalizer with
target $C$ exists if and only if $(0,1)\in L$. Suppose not, and let $x\ne0$
be the class of $(0,1)$ in $F/L$. There is a homomorphism
$\psi\colon F/L\to\R/\Z$ such that $c=\psi(x)$ has distance at least $1/3$
from $0$: on the cyclic group generated by $x$ send $x$ to $1/2$ if $x$ has
infinite order and to $\lfloor m/2\rfloor/m$ if $x$ has order $m\ge2$, and
extend to $F/L$, which is possible because $\R/\Z$ is divisible. The
homomorphism $\psi$ is given by functions
$\theta_1,\dots,\theta_k\colon\C(P,C)\to\R/\Z$ and the value $c$ on $(0,1)$,
and it vanishes on $L$:
\[
 \sum_{i=1}^k\bigl(\theta_i(ga_i)-\theta_i(ga_0)\bigr)+c=0
 \qquad\text{for every }g\in\C(Q,C).
\]
Divide $\R/\Z$ into $N$ arcs of length $1/N$ and colour $e\in\C(P,C)$ by the
$k$-tuple of the arcs containing $\theta_1(e),\dots,\theta_k(e)$. For a $g$
as in the Ramsey property each difference $\theta_i(ga_i)-\theta_i(ga_0)$
has distance less than $1/N$ from $0$, so $c$ has distance less than
$k/N<1/3$ from $0$, a contradiction.

\emph{Property \textup{(Re)}.} Let $N>2k/\varepsilon$ and let $C$ be as in
Definition~\ref{def:RP} for $r=N^k$. Let $\Delta$ be the simplex of
probability vectors on $\C(Q,C)$ and
$\Lambda=([-1,1]^{\C(P,C)})^k$. The function
$\Phi(\nu,\lambda)=\sum_i\langle\lambda_i,\nu a_i-\nu a_0\rangle$ is bilinear
on the compact convex sets $\Delta$ and $\Lambda$, so by the minimax theorem \cite{vN} (applied to the matrix game whose columns are the vertices of the cube $\Lambda$; a linear function on $\Lambda$ takes its maximum at a vertex) $\min_\nu\max_\lambda\Phi=\max_\lambda\min_\nu\Phi$. Here
$\max_\lambda\Phi(\nu,\lambda)=\sum_i\|\nu a_i-\nu a_0\|_1$, and
$\min_\nu\Phi(\nu,\lambda)$ is the minimum over $g\in\C(Q,C)$ of
$\sum_i(\lambda_i(ga_i)-\lambda_i(ga_0))$. Given $\lambda$, colour
$e\in\C(P,C)$ by the $k$-tuple of the intervals of length $2/N$ of $[-1,1]$
containing $\lambda_i(e)$; a $g$ as in the Ramsey property gives a sum of
absolute value at most $2k/N<\varepsilon$. So the common value is less than
$\varepsilon$.
\end{proof}

\begin{remark}\label{rem:evans}
For the category of finite substructures of the Fra\"iss\'e limit $M^*$ of a Fra\"iss\'e class of finite structures that have a linear order among their relations, the Ramsey property is equivalent to extreme amenability of $\Aut(M^*)$, by the theorem of Kechris, Pestov and Todorcevic \cite[Theorem~4.8]{KPT}, which is stated for such order classes and for colourings of substructures; for ordered structures these are the embeddings; and for an extremely amenable group Evans proves
that invariants are exact on compact continuous modules
\cite[Lemma~3.7\,(i)]{Ev}, which by Pontryagin duality is the statement that
coinvariants are exact on smooth $\Z[G]$-modules. With
Proposition~\ref{prop:affine} this gives the first assertion of
Lemma~\ref{lem:ramsey} for such categories. The finite proof above, which
comes from the runs, is included because Section~\ref{sec:vector} applies
the lemma to a category for which we want no limit object, and because it
is short.
\end{remark}

\begin{lemma}\label{lem:Re}
Assume \textup{(K1)}. If $\A(M,G,\K)$ has property \textup{(Re)}, then $G$
is amenable.
\end{lemma}

\begin{proof}
Let $G$ act continuously on a non-empty compact Hausdorff space $X$. The
invariant probability measures are the intersection, in the compact space
of all regular Borel probability measures with the weak* topology, of the closed
sets $\{\lambda:|\lambda(f\circ g)-\lambda(f)|\le\eta\}$ for continuous
$f\colon X\to[0,1]$, $g\in G$ and $\eta>0$. So it is enough to find, for
continuous $f_1,\dots,f_r\colon X\to[0,1]$, elements $g_1,\dots,g_s\in G$ and
$\eta>0$, a probability measure $\lambda$ with
$|\lambda(f_i\circ g_j)-\lambda(f_i)|\le3\eta$ for all $i,j$.

By continuity of the action and compactness of $X$ there is $P\in\K$ with
$|f_i(ux)-f_i(x)|<\eta$ for all $x\in X$, $u\in G_P$ and all $i$. Fix
$x_0\in X$. For each map $e=g|_P$, $g\in G$, choose one such $g$ and put
$c_i(e)=f_i(g^{-1}x_0)$; for any other $g'$ with $g'|_P=e$ we have
$g'=gu$ with $u\in G_P$, so $|f_i(g'^{-1}x_0)-c_i(e)|<\eta$. Let $Q\in\K$
contain $P$ and all $g_j^{-1}P$, let $a_0\colon P\to Q$ be the inclusion and
$a_j=g_j^{-1}|_P$. By \textup{(Re)} there are $C$ and a probability vector
$\nu$ on $\A(Q,C)$ with $\sum_j\|\nu a_j-\nu a_0\|_1<\eta$. Choose
$\hat h\in G$ extending each $h\in\A(Q,C)$ and let
$\lambda=\sum_h\nu(h)\delta_{\hat h^{-1}x_0}$. Since $\hat h$ extends $ha_0$
and $\hat hg_j^{-1}$ extends $ha_j$,
\[
 \bigl|\lambda(f_i)-\textstyle\sum_h\nu(h)c_i(ha_0)\bigr|<\eta,\qquad
 \bigl|\lambda(f_i\circ g_j)-\textstyle\sum_h\nu(h)c_i(ha_j)\bigr|<\eta ,
\]
and $|\sum_h\nu(h)(c_i(ha_j)-c_i(ha_0))|=|\langle c_i,\nu a_j-\nu
a_0\rangle|\le\|\nu a_j-\nu a_0\|_1<\eta$.
\end{proof}

\subsection{Measures on expansions and the transfer theorem}

Let $H$ be a subgroup of $G$ and $E(B)=G_B\backslash G/H$ for $B\in\K$. For
$B\subseteq C$ there is the surjection $E(C)\to E(B)$,
$\zeta\mapsto\zeta|_B$, induced by $G_C\subseteq G_B$, and for $a\in G$ the
bijection $a_*\colon E(B)\to E(aB)$, $G_BgH\mapsto G_{aB}agH$.

\begin{definition}\label{def:measure}
A \emph{normalized invariant measure} with values in a ring $k$ is a family
of functions $\mu_B\colon E(B)\to k$, $B\in\K$, each $E(B)$ being finite,
such that
\begin{itemize}
\item[\textup{(i)}] $\sum_{\xi\in E(B)}\mu_B(\xi)=1$;
\item[\textup{(ii)}] $\mu_B(\xi)=\sum_{\zeta|_B=\xi}\mu_C(\zeta)$ for
$B\subseteq C$ and $\xi\in E(B)$;
\item[\textup{(iii)}] $\mu_{aB}(a_*\xi)=\mu_B(\xi)$ for $a\in G$.
\end{itemize}
\end{definition}

Equivalently, it is a $G$-invariant finitely additive function of total
mass $1$ on the closed and open subsets of the compact space
$X=\varprojlim_BE(B)$, on which $G$ acts continuously by
$(ax)_{aB}=a_*(x_B)$.

\begin{proof}[Proof of Theorem~\ref{thm:transfer}]
(a) Let $R$ be a $k$-algebra and $N$ a smooth $R[G]$-module. For $n\in N$
choose $B\in\K$ with $G_B$ fixing $n$, and for $\xi=G_BgH\in E(B)$ let
$f_n(\xi)$ be the class of $g^{-1}n$ in $N_H$. It is well defined:
$(ugh)^{-1}n=h^{-1}g^{-1}n$ for $u\in G_B$, $h\in H$, and this has the same
class as $g^{-1}n$. Put
\[
 s_N(n)=\sum_{\xi\in E(B)}\mu_B(\xi)f_n(\xi)\in N_H .
\]
By (ii) this does not depend on $B$: for $B\subseteq C$ the function $f_n$
on $E(C)$ is the composite of $E(C)\to E(B)$ with $f_n$ on $E(B)$, and any
two choices are contained in a third. So $s_N$ is $R$-linear, and it is
natural in $N$. For $a\in G$ the element $an$ is fixed by $G_{aB}$ and
$f_{an}(a_*\xi)=f_n(\xi)$, so $s_N(an)=s_N(n)$ by (iii). Hence $s_N$ induces
$\bar s_N\colon N_G\to N_H$. If $q_N\colon N_H\to N_G$ is the quotient map,
then $q_N(f_n(\xi))$ is the class of $n$, and $q_N\bar s_N=\id$ by (i). Now
let $N'\subseteq N$ be smooth and let $x\in N'_G$ go to zero in $N_G$. Then
$\bar s_{N'}(x)$ goes to zero in $N_H$; as $N'_H\to N_H$ is injective for
$R[H]$-modules that are smooth for $H$ (the hypothesis on $H$ over $k$
gives it over $R$, since coinvariants do not depend on the ring),
$\bar s_{N'}(x)=0$ and $x=q_{N'}\bar s_{N'}(x)=0$. By Theorem~\ref{thm:dict},
$\A$ is exact over $R$.

(b) Let $F$ be the $k$-vector space of the functions $G/H\to k$ that are
invariant under $G_B$ for some $B\in\K$, with $(a\varphi)(x)=\varphi(a^{-1}x)$.
It is a smooth $k[G]$-module containing the constants. By exactness and
Theorem~\ref{thm:dict}, $k=k_G\to F_G$ is injective, so there is a linear
form $\ell$ on $F_G$ with value $1$ on the class of the constant function
$1$. Let $\mu_B(\xi)$ be the value of $\ell$ on the class of the indicator
function $1_\xi$ of $\xi\subseteq G/H$. Then (i) and (ii) hold because
$\sum_\xi1_\xi=1$ and $1_\xi=\sum_{\zeta|_B=\xi}1_\zeta$, and (iii) because
$a1_\xi=1_{a_*\xi}$.

(c) The space $X$ is non-empty and compact, so there is a $G$-invariant
Borel probability measure $\lambda$ on it. The measures of the sets
$\{x:x_B=\xi\}$ satisfy (i)--(iii) with real values. Composing with a
$\Q$-linear map $\R\to\Q$ that sends $1$ to $1$ gives a normalized invariant
measure with values in $\Q$, hence in $k$, and by (a) the category $\A$ is exact over $k$. By Proposition~\ref{prop:affine} every finite family of arrows then has an affine equalizer $t$ over $k$. If $\rho\colon k\to\Q$ is $\Q$-linear with $\rho(1)=1$, applying $\rho$ to the coefficients of $t$ gives an affine equalizer over $\Q$, because the conditions $\aug(t)=1$ and $ta_i=ta_0$ are linear with integer coefficients. So $\A$ is exact over $\Q$, hence over every $\Q$-algebra.
\end{proof}

\begin{remark}\label{rem:zucker}
Let $H=\Aut(M^*)$ for an expansion $M^*$ of a homogeneous structure $M$, as
in Section~\ref{sec:posets}. Zucker \cite[\S3]{Zu} considers the rational
vector space on the set $\Omega$ of expansions of finite substructures and
its subspace $V$ spanned by the vectors ``an expansion minus the sum of its
extensions along an embedding''. His Theorem~3.1 states, for an excellent
pair, that $\Aut(M)$ is amenable if and only if every non-zero vector of $V$
has coefficients of both signs. A normalized invariant measure with values
in a field $k$ is a linear form on $k\Omega$ that vanishes on $V$ and takes
the value $1$ on the expansion of the empty structure. So, for $H$ as in
Theorem~\ref{thm:transfer}, the category is exact over $k$ if and only if
that vector does not lie in $V\otimes k$.
\end{remark}

\section{Posets with an independent linear order}\label{sec:posets}

\subsection{Integer weights on linear extensions}

Let $(P,\prec,<)$ be in $\Po$; we call $<$ the reference order. A
\emph{cell} of $P$ is a sequence $(B_1,\dots,B_b)$ of non-empty pairwise
disjoint subsets with union $P$ such that $y\prec z$ implies that the block
of $y$ comes before the block of $z$. A cell \emph{sorts to} the linear
order on $P$ that lists $B_1,\dots,B_b$ in this order and each block in the
reference order; this is a linear extension of $\prec$. For a linear
extension $\pi$ of $\prec$ put
\[
 w_P(\pi)=(-1)^{|P|}\sum_{\text{cells sorting to }\pi}(-1)^{b},
\]
where $b$ is the number of blocks. The empty structure has one cell, with
$b=0$, and $w_\varnothing=1$.

\begin{lemma}\label{lem:weights}
The integers $w_P(\pi)$ are invariant under isomorphisms of structures
$(P,\prec,<)$, their sum over the linear extensions of $P$ is $1$, and for
every subset $Q\subseteq P$, with the induced relations, and every linear
extension $\tau$ of $Q$,
\[
 \sum_{\pi|_Q=\tau}w_P(\pi)=w_Q(\tau).
\]
\end{lemma}

\begin{proof}
Invariance is clear. By induction it is enough to prove the identity for
$P=Q\cup\{x\}$; the case $Q=\varnothing$ of the general identity is the
statement on the sum. Deleting $x$ from a cell of $P$, and its block if it
becomes empty, gives a cell of $Q$, and the sorting of the new cell is the
restriction of the sorting of the old one. Fix a cell $(B_1,\dots,B_b)$ of
$Q$. Let $l$ be the largest index of a block containing an element below
$x$ in $\prec$, or $0$, and $u$ the smallest index of a block containing an
element above $x$, or $b+1$. Then $l<u$, since $y\prec x\prec z$ implies
$y\prec z$. The cells of $P$ over the given one are obtained by adding $x$
to a block $B_j$ with $l<j<u$, which gives $m=u-l-1$ cells with $b$ blocks,
or by inserting $\{x\}$ as a new block after $B_j$ for some $l\le j<u$,
which gives $m+1$ cells with $b+1$ blocks. Their signed number is
$m(-1)^b+(m+1)(-1)^{b+1}=(-1)^{b+1}$, and
$(-1)^{|P|}(-1)^{b+1}=(-1)^{|Q|}(-1)^b$. Summing over the cells of $Q$ that
sort to $\tau$ gives the identity.
\end{proof}

\begin{remark}\label{rem:HS}
The weights are values of a measure of Harman and Snowden. Let $O(P)$ be
the set of maps $t\colon P\to\R$ with $t_y<t_z$ for $y\prec z$, an open
convex cone in $\R^P$, and $O_\pi(P)$ the set of those $t$ for which
ordering $P$ by the values of $t$, ties broken by the reference order,
gives $\pi$. The cells above are the orbits of $\Aut(\R,<)$ on $O(P)$; a
cell with $b$ blocks is the set of strictly increasing $b$-tuples of reals.
The compactly supported Euler characteristic is a $\Z$-valued measure
$\mu$ for $\Aut(\R,<)$ in the sense of \cite{HS} (Example~3.5,
Theorem~16.7 and Remark~16.8 there), with value $(-1)^b$ on such a cell. So
\[
 w_P(\pi)=\mu(O_\pi(P))/\mu(O(P)),\qquad\mu(O(P))=(-1)^{|P|},
\]
and the identity of Lemma~\ref{lem:weights} is the transitivity of
push-forward \cite[Proposition~4.8]{HS} for the projections $O(P)\to O(P\setminus\{x\})$, whose fibres are open intervals. The family
$(w_P)$ is not itself a measure in the sense of \cite{HS} on the class of
expanded structures, because it is not multiplicative along embeddings: for
the antichain $\{x,y\}$ with $x<y$ one finds $w(x,y)=0$ and $w(y,x)=1$,
while after adding a point $z$ with $x\prec z$, incomparable to $y$, and
$y<z$, the three linear extensions have weights $w(x,y,z)=-1$,
$w(x,z,y)=1$, $w(y,x,z)=1$.
\end{remark}

\subsection{Proof of Theorem~\ref{thm:posets}}

Let $\Po^*$ be the class of finite structures $(P,\prec,<,\lhd)$ with
$(P,\prec,<)\in\Po$ and $\lhd$ a linear extension of $\prec$. Both $\Po$
and $\Po^*$ are Fra\"iss\'e classes \cite[Remark~3.2]{KS}, and $\Po^*$ has
the Ramsey property for embeddings (Bodirsky \cite[Corollary~1.4\,(3)]{Bo}).
Let $M^*=(M,\prec,<,\lhd_0)$ be the Fra\"iss\'e limit of $\Po^*$. Every
linear extension of a subposet of a finite poset extends to the poset, so
$(M,\prec,<)$ has the extension property that characterizes the Fra\"iss\'e
limit of $\Po$; we take $M_{\Po}=(M,\prec,<)$. Let $G=\Aut(M_{\Po})$,
$H=\Aut(M^*)\subseteq G$ and $\K$ the set of finite subsets of $M$.

\emph{The group $H$.} The category $\A(M,H,\K)$ is the category of
non-empty finite substructures of $M^*$ and embeddings; it has the Ramsey
property because $\Po^*$ has it and every member of $\Po^*$ embeds in
$M^*$. By Lemma~\ref{lem:ramsey} and Proposition~\ref{prop:affine} it is
exact over $\Z$, so coinvariants preserve monomorphisms of smooth
$R[H]$-modules for every ring $R$ (Theorem~\ref{thm:dict}).

\emph{The sets $E(B)$.} For $B\in\K$ the map $gH\mapsto(g\lhd_0)|_B$
induces a bijection from $E(B)=G_B\backslash G/H$ to the set of linear
extensions of $(B,\prec)$; here $g\lhd_0$ is the image of $\lhd_0$ under
$g$. Indeed, if $(g\lhd_0)|_B=(g'\lhd_0)|_B$, then $g'^{-1}g$ restricts to
an isomorphism of the substructures $g^{-1}B$ and $g'^{-1}B$ of $M^*$, which
extends to some $h\in H$, and $g'hg^{-1}\in G_B$. And every linear extension
$\pi$ occurs: $(B,\prec,<,\pi)$ embeds in $M^*$ by a map $e$, and by
homogeneity of $M_{\Po}$ there is $g\in G$ with $g^{-1}|_B=e$. In
particular $E(B)$ is finite. Under these bijections the maps
$E(C)\to E(B)$ and $a_*$ are restriction and transport of linear orders.

\emph{Exactness.} By Lemma~\ref{lem:weights} the functions $\mu_B=w_B$ form
a normalized invariant measure with values in $\Z$. By
Theorem~\ref{thm:transfer}\,(a) the category $\A(M,G,\K)$, which is
equivalent to the category of non-empty members of $\Po$, is exact over
every ring.

\emph{Non-amenability} is \cite[Remark~3.2]{KS}: the class called
$\mathcal{OP}$ there is $\Po$. \qed

\begin{proof}[Proof of Corollary~\ref{cor:evans}]
By Theorems~\ref{thm:posets} and~\ref{thm:dict} and Lemma~\ref{lem:open},
$N^U\to N_U$ is injective for every smooth $R[U]$-module $N$. Apply this to
$N=R[W]/Y$: the map $\Omega_U$ identifies $R[W]_U$ with $R[U\backslash W]$,
so $N_U=R[U\backslash W]/\Omega_U(Y)$. If $\Omega_U(x)\in\Omega_U(Y)$, the
class of $x$ in $N$ is $U$-invariant and goes to zero in $N_U$, hence is
zero, and $x\in Y$. If $R$ is a field and $x\notin Y$, take a linear form on
$R[U\backslash W]$ that vanishes on $\Omega_U(Y)$ and not on $\Omega_U(x)$,
and compose it with $\Omega_U$.
\end{proof}

\begin{remark}\label{rem:reducts}
The proof uses only that the category is exact over $R$. So the conclusions
of Corollary~\ref{cor:evans} hold over $R$ whenever $\A(M,G,\K)$ is exact
over $R$; by Theorem~\ref{thm:transfer}\,(c) and Remark~\ref{rem:evans} this
is the case over every field of characteristic $0$ when $M$ is a reduct of
an $\omega$-categorical Ramsey structure $M^*$ and $\Aut(M)$ is amenable.
Conversely, the first conclusion of Corollary~\ref{cor:evans}, required
only for $U=G$, implies that $\A(M,G,\K)$ is exact over $R$. Let
$N'\subseteq N$ be smooth and $n'\in N'\cap I_GN$. The class of $n'$ in
$N/I_GN'$ is $G$-invariant. Write $N/I_GN'=R[W]/Y$, with $W$ a $G$-set with
open stabilizers and that class the image of a point $w$ of $W$ fixed by
$G$. The class goes to zero in $(N/I_GN')_G=N_G$, so
$\Omega_G(w)\in\Omega_G(Y)$; hence $w\in Y$ and $n'\in I_GN'$. So by
Theorem~\ref{thm:vector}\,(c) the membership test with $U=G$ fails for the
ordered $\F_2$-space over every ring in which $2$ is not invertible. This
does not answer the question of \cite[Remarks~2.12]{Ev}, which asks only for
a function invariant under the pointwise stabilizer of some finite set: for
a finite group acting faithfully that weaker property holds trivially, while
coinvariants are not exact in the modular case.
\end{remark}

\section{Ordered vector spaces over the field with two
elements}\label{sec:vector}

The objects of $\V$ are the pairs $(W,<)$ with $W$ a finite-dimensional
non-zero $\F_2$-vector space and $<$ a linear order on $W\setminus\{0\}$; the
arrows are the injective linear maps that preserve the order. Two objects
embed in a third, and two arrows $W_0\to W_1$, $W_0\to W_2$ can be completed
to a commutative square in which $W_1\cap W_2=W_0$: take the push-out of
vector spaces, an order on the union of $W_1$ and $W_2$ that extends both
orders, and the remaining vectors last. There are countably many objects up
to isomorphism. By Fra\"iss\'e's theorem for classes of finitely generated
structures (\cite[\S7.1]{Ho}; see also \cite[Theorem~2.2]{KPT}), applied to the language with addition and the
order, there is a countable vector space $V_\infty$ with a linear order on
$M=V_\infty\setminus\{0\}$ such that every object of $\V$ embeds in
$(V_\infty,<)$ and every isomorphism in $\V$ between finite-dimensional
subspaces of $V_\infty$ extends to an element of the group $G$ of
order-preserving linear automorphisms. Let $\K$ consist of $\varnothing$
and the sets $W\setminus\{0\}$ for finite-dimensional subspaces $W\ne0$.
Then \textup{(K1)} holds and $\A(M,G,\K)$ is equivalent to $\V$.

\subsection{Amenability}

\begin{proof}[Proof of Theorem~\ref{thm:vector}\,\textup{(a)}]
By Lemma~\ref{lem:Re} it is enough to prove that the category $\V$ has
property \textup{(Re)}. Let $a_0,\dots,a_t$ be arrows from $P$ to $Q$ in
$\V$ and let $\varepsilon>0$. Let $c$ and $d$
be the dimensions of $P$ and $Q$, and $p=2^c-1$, $q=2^d-1$ the numbers of
their non-zero vectors. The orders identify the non-zero vectors of $P$
with $\{1,\dots,p\}$ and those of $Q$ with $\{1,\dots,q\}$, and they
identify $a_i$ with an increasing injection $\alpha_i$ of
$\{1,\dots,p\}$ into $\{1,\dots,q\}$.

\emph{Step 1.} The category of finite linear orders and increasing
injections has the Ramsey property, by Ramsey's theorem, hence
\textup{(Re)} by Lemma~\ref{lem:ramsey}. So there are $m$ and a probability
vector $\nu$ on the increasing injections
$\{1,\dots,q\}\to\{1,\dots,m\}$ with
$\sum_{i=1}^t\|\nu_i-\nu_0\|_1<\varepsilon/3$, where $\nu_i$ is the image
of $\nu$ under composition with $\alpha_i$. Extend $\nu$ and the $\nu_i$ by
zero to all maps into $\{1,\dots,m\}$.

\emph{Step 2.} Let $n\ge d$ and $V=\F_2^n$. Colour the non-zero vectors of
$V$ independently and uniformly by $\chi(v)\in\{1,\dots,m\}$, and order
$V\setminus\{0\}$ by colour, ties broken by a fixed order. For an injective
linear map $g\colon Q\to V$ put $w(g)=\nu(\chi\circ g)$; it is non-zero only
if $g$ preserves the orders. Let $e_r(n)=\prod_{j<r}(2^n-2^j)$ be the
number of injective linear maps $\F_2^r\to V$, and $W=\sum_gw(g)$. The
colours of the $q$ points $g(Q\setminus\{0\})$ are independent and uniform,
so the expectation of $w(g)$ is $m^{-q}$ and
\[
 \mathbb{E}\,W=e_d(n)\,m^{-q}.
\]
For an injective linear $u\colon P\to V$ let
$X_i(u)=\sum_{ga_i=u}w(g)$; there are $N=e_d(n)/e_c(n)$ such $g$. Given the
colours on $u(P)$, the colours of the $q-p$ other points of $g(Q)$ are
independent and uniform, so the conditional expectation of $X_i(u)$ is
\[
 H_i(u)=N\,m^{-(q-p)}\,\nu_i(\chi\circ u).
\]
Two maps $g,h$ with $ga_i=ha_i=u$ and $g(Q)\cap h(Q)=u(P)$ have
conditionally independent weights. For fixed $g$ and $h$ uniform among the
$N$ maps, each vector of $Q$ outside $a_i(P)$ has an image that is uniform
on $V\setminus u(P)$, so the probability that $g(Q)\cap h(Q)\ne u(P)$ is at
most $\eta_n=(2^d-2^c)^2/(2^n-2^c)$. As $0\le w\le1$, the conditional
variance of $X_i(u)$ is at most $N^2\eta_n$, so
$\mathbb{E}|X_i(u)-H_i(u)|\le N\sqrt{\eta_n}$, and summing over the
$e_c(n)$ maps $u$,
\[
 \mathbb{E}\|X_i-H_i\|_1\le e_d(n)\sqrt{\eta_n},\qquad
 \mathbb{E}\|H_i-H_0\|_1=e_d(n)\,m^{-q}\,\|\nu_i-\nu_0\|_1 ,
\]
the second because the $p$ colours on $u(P)$ are independent and uniform.
Let $D=\sum_{i=1}^t\|X_i-X_0\|_1$. Then
\[
 \mathbb{E}\,D\le e_d(n)\Bigl(2t\sqrt{\eta_n}
 +m^{-q}\sum_{i=1}^t\|\nu_i-\nu_0\|_1\Bigr).
\]
Choose $n$ with $2tm^q\sqrt{\eta_n}<\varepsilon/3$. Then
$\mathbb{E}(D-\varepsilon W)<0$, so for some colouring $D<\varepsilon W$,
and $W>0$. For this colouring $(V,<)$ is an object of $\V$, $w/W$ is a
probability vector on $\V(Q,V)$, its images under the $a_i$ are the
$X_i/W$, and $\sum_i\|X_i/W-X_0/W\|_1<\varepsilon$.
\end{proof}

\subsection{The obstruction in characteristic two}

\begin{proof}[Proof of Theorem~\ref{thm:vector}\,\textup{(c)}]
Let $P=\F_2^2$ with non-zero vectors $x_1<x_2<x_3$. Write the non-zero
vectors of $\F_2^3$ as $1,\dots,7$ in binary, so that addition is bitwise
addition without carry, and let $Q$ be $\F_2^3$ with the order
$1<2<3<4<6<5<7$. Let $a_1,\dots,a_5\colon P\to Q$ send $(x_1,x_2,x_3)$ to
\[
 (1,2,3),\quad(1,6,7),\quad(2,5,7),\quad(3,4,7),\quad(3,6,5).
\]
Each triple is the set of non-zero vectors of a plane, listed increasingly,
and every bijection between such sets is linear; so these are arrows of
$\V$.

Let $D$ be an object of $\V$. Choose a basis $b_1,\dots,b_n$ of $D$ and let
$\sqsubset$ be the order on $D$ that compares two vectors at the largest
index where their coordinates differ, with $0$ before $1$; it has nothing to
do with the given order of $D$. For a statement $S$ let $[S]$ be $1$ if $S$
holds and $0$ otherwise. For $v\ne0$ let $\epsilon_v$ be the linear
form that takes the coordinate at the largest index where $v$ has
coordinate $1$. For $x\ne y$ the vectors $x$ and $y$ differ last at that
index of $x+y$, so, with values in $\F_2$,
\begin{equation}\label{eq:ind}
 [x\sqsubset y]=\epsilon_{x+y}(y).
\end{equation}
For distinct non-zero $u,v$,
\begin{equation}\label{eq:between}
 [u+v\text{ lies between }u\text{ and }v]=\epsilon_u(v)+\epsilon_v(u).
\end{equation}
Indeed, if the largest indices of $u$ and $v$ agree, then $u+v$ comes before
both and the right side is $1+1$. If the largest index of $u$ exceeds that
of $v$, then $v$ comes before $u$ and before $u+v$, the vector $u+v$ comes
before $u$ exactly when $\epsilon_v(u)=1$, and $\epsilon_u(v)=0$.

Let $g\colon Q\to D$ be linear and injective and write $i$ for $g(i)$. By
\eqref{eq:ind}, and since $1+2=5+6=3$ and $1+6=2+5=7$,
\[
 [2\sqsubset1]+[1\sqsubset6]+[2\sqsubset5]+[5\sqsubset6]
 =\epsilon_3(1)+\epsilon_7(6)+\epsilon_7(5)+\epsilon_3(6)
 =\epsilon_3(7)+\epsilon_7(3),
\]
which by \eqref{eq:between} is $[4\text{ lies between }3\text{ and }7]$. For
an arrow $e\colon P\to D$ define, with $y_j=e(x_j)$,
\begin{gather*}
 f_1(e)=[y_2\sqsubset y_1],\qquad f_2(e)=f_3(e)=[y_1\sqsubset y_2],\\
 f_4(e)=[y_2\text{ lies between }y_1\text{ and }y_3],\qquad
 f_5(e)=[y_3\sqsubset y_2].
\end{gather*}
The computation says $\sum_jf_j(ga_j)=0$ for every arrow $g\colon Q\to D$.
On the other hand $\sum_jf_j(e)=1$ for every $e$: with
$\alpha=[y_2\sqsubset y_1]$ and $\beta=[y_3\sqsubset y_2]$ one has
$f_4=(1-\alpha)(1-\beta)+\alpha\beta=1+\alpha+\beta$ in $\F_2$.

Let $R$ be a ring in which $2$ is not invertible, and suppose
$t\in R[\V(Q,D)]$ is an affine equalizer of $a_1,\dots,a_5$. Reduce modulo
$2$: in the non-zero $\F_2$-algebra $\bar R=R/2R$ we get $\bar t$ with
$\aug\bar t=1$ and $\bar ta_j=\bar m$ for all $j$. Extending the $f_j$
linearly,
\[
 0=\sum_jf_j(\bar ta_j)=\sum_jf_j(\bar m)=\aug(\bar m)=\aug(\bar t)=1 ,
\]
a contradiction. By Proposition~\ref{prop:affine}, applied to
$\A(M,G,\K)\simeq\V$, the category is not exact over $R$.
\end{proof}

\subsection{Exactness over the rationals}

Let $\C$ be the following category. An object is a finite non-empty linearly
ordered set $E$ with a linear order on the set of its non-empty subsets; an
arrow is an increasing injection that preserves the order on subsets. The
functor $u\colon\C\to\V$ sends $E$ to the vector space $\F_2^E$, whose
non-zero vectors are the sums of non-empty sets of basis vectors, with the
corresponding order. For a vector of $\F_2^E$, its \emph{support} is the set
of basis vectors that occur in it.

\begin{lemma}\label{lem:C}
The category $\C$ has the Ramsey property, and it is exact over $\Z$.
\end{lemma}

\begin{proof}
Let $A,B$ be objects and $b=|B|$; we may assume $|A|\le b$. For a set $S$
of size at most $b$ in a linearly ordered set, with elements
$s_1<\dots<s_j$, let $\tau(S)=(s_1,\dots,s_j,s_j,\dots,s_j)$, a $b$-tuple.
Regard $A$ and $B$ as structures $A',B'$ with a linear order on the points,
the $b$-ary relation consisting of the tuples $\tau(S)$, and the linear
order on this relation given by the order on subsets. By the theorem of
Hubi\v cka and Ne\v set\v ril on structures with linearly ordered points and
linearly ordered relations \cite[Theorem~4.46]{HN} there is, for every $r$,
such a structure $C_0$ with the Ramsey property for $A'$, $B'$ and $r$
colours, embeddings preserving the relation, its complement and both
orders. On the points of $C_0$ define an object $C$ of $\C$: order the
subsets of size at most $b$ so that the order agrees with that of $C_0$ on
those $S$ whose tuple $\tau(S)$ is in the relation of $C_0$, and put the
other subsets anywhere. An embedding $A'\to C_0$ or $B'\to C_0$ maps the
tuples $\tau(S)$ to tuples of the same form in the relation of $C_0$, hence
is an arrow of $\C$ into $C$; and $\C(A,B)$ is the set of embeddings
$A'\to B'$. So a colouring of $\C(A,C)$ restricts to a colouring of the
embeddings $A'\to C_0$, and an embedding $g\colon B'\to C_0$ with all
$g\circ\C(A,B)$ of one colour is the required arrow. This proves the Ramsey
property.

The objects of $\C$ are the finite structures in a relational language,
with relations saying that one increasing tuple precedes another as a
subset, and the arrows are the embeddings. This class is hereditary, has
countably many members up to isomorphism, and has the amalgamation and
joint embedding properties: on the union of two members over a common part,
merge the two orders of the points, extend the two orders on subsets to a
common order on the subsets of either member, and put the remaining
subsets last. By Fra\"iss\'e's theorem $\C$ is equivalent to a category
$\A(M',G',\K')$ as in Section~\ref{sec:intro}, and
Lemma~\ref{lem:ramsey} and Proposition~\ref{prop:affine} give exactness
over $\Z$.
\end{proof}

\begin{proof}[Proof of Theorem~\ref{thm:vector}\,\textup{(b)}]
Let $R$ be a $\Q$-algebra and $a_0,\dots,a_k\colon P\to Q$ arrows of $\V$.
By Proposition~\ref{prop:affine} it is enough to find an affine equalizer
over $\Q$. For an object $Y$ of $\V$ let $X_Y\colon\C\to\Q\text{-Mod}$ be
the functor $C\mapsto\Q[\V(Y,uC)]$; write $[C,h]$ for the class of
$h\in\V(Y,uC)$ in $\colim_\C X_Y$.

\emph{Step 1: classes in $\colim_\C X_P$.} For $h\in\V(P,uC)$ let
$U_h\subseteq C$ be the union of the supports of the vectors of $h(P)$,
with the induced orders; it is an object of $\C$, the inclusion
$U_h\subseteq C$ is an arrow, and $h$ factors through $u(U_h)$. So
$[C,h]=[U_h,h]$, and $[U,h]=[U',h']$ whenever there is an isomorphism
$\varphi\colon U\to U'$ in $\C$ with $u(\varphi)h=h'$. Call the isomorphism
class of the pair $(U_h,h)$ the \emph{type} of $(C,h)$. Then a finite
combination $\sum_jr_j[C_j,h_j]$ depends only on the sums of the
coefficients over each type.

\emph{Step 2: a class in $\colim_\C X_Q$ with equal images.} Let $Q^\vee$
be the dual space, $E=Q^\vee\setminus\{0\}$, and
$j\colon Q\to\F_2^E$, $j(v)=(\ell(v))_{\ell\in E}$, which is linear and
injective. Let $\Omega$ be the finite set of pairs $\omega$ consisting of a
linear order on $E$ and a linear order on the non-zero vectors of $\F_2^E$
in which the vectors outside $j(Q)$ come first, in any order, followed by
$j(Q\setminus\{0\})$ in the order of $Q$. Each $\omega$ gives an object
$E_\omega$ of $\C$ and $j\in\V(Q,uE_\omega)$. Put
\[
 \mu=\frac1{|\Omega|}\sum_{\omega\in\Omega}[E_\omega,j]\in\colim_\C X_Q,
 \qquad\aug\mu=1 .
\]
The image of $\mu$ under $a_i$ is the average of the classes
$[E_\omega,ja_i]$. The union of the supports of the vectors $ja_i(x)$ is
$E_i=\{\ell\in E:\ell a_i\ne0\}$, and
\begin{equation}\label{eq:inter}
 \F_2^{E_i}\cap j(Q)=ja_i(P):
\end{equation}
$j(v)$ has support in $E_i$ exactly when every linear form vanishing on
$a_i(P)$ vanishes on $v$, that is, when $v\in a_i(P)$. So the type of
$(E_\omega,ja_i)$ is that of the set $E_i$ with the order induced by
$\omega$, the order on the non-zero vectors of $\F_2^{E_i}$ in which those
outside $ja_i(P)$ come first, in the order induced by $\omega$, followed by
$ja_i(P\setminus\{0\})$ in the order of $P$, and the map $ja_i$. As
$\omega$ runs over $\Omega$, the induced order on $E_i$ and the induced
order on the vectors of $\F_2^{E_i}$ outside $ja_i(P)$ are independent and
uniformly distributed, because restriction of linear orders to a subset
has fibres of equal size. Finally, for two indices $i,i'$ there is a
bijection $\beta\colon E_i\to E_{i'}$ with $\beta(\ell)a_{i'}=\ell a_i$,
because every non-zero linear form on $P$ is $\ell a_i$ for exactly
$2^{\dim Q-\dim P}$ elements $\ell$ of $E_i$, and likewise for $i'$; the
linear map defined by $\beta$ carries $ja_i$ to $ja_{i'}$. Hence the
distribution of the type of $(E_\omega,ja_i)$ is the same for all $i$, and
by Step~1 the images of $\mu$ under $a_0,\dots,a_k$ coincide.

\emph{Step 3: an affine equalizer.} Let $D\colon X_Q\to X_P^k$ be the
natural transformation with components $t\mapsto ta_i-ta_0$. By
Lemma~\ref{lem:C} the functor $\colim_\C$ is exact, so it commutes with
kernels, and $\mu$, which lies in the kernel of $\colim_\C D$, is the class
of an element of $\colim_\C\ker D$. Such an element is a finite sum of
classes of elements of $(\ker D)(C_j)$; any two objects of $\C$ map to a
third, so it is the class of a single $t\in(\ker D)(C)\subseteq
\Q[\V(Q,uC)]$. The augmentation is natural, so $\aug t=\aug\mu=1$, and
$D(C)t=0$ says $ta_i=ta_0$.
\end{proof}

\section{Provenance and checks}\label{sec:provenance}

This paper is AI-produced. No human has proved, written or checked any part
of it.

\paragraph{What is whose.} The statements and proofs of
Theorems~\ref{thm:posets} and~\ref{thm:vector} (part~(c) over $\F_2$ and
$\Z$), of the first sentence of
Theorem~\ref{thm:dict}, of Theorem~\ref{thm:transfer} in the form of a
group with a closed co-precompact subgroup and a finitely additive
invariant function on the closed and open sets of the completion of $G/H$,
of Lemmas~\ref{lem:Re}, \ref{lem:weights} and~\ref{lem:C} and of the three
steps of Section~\ref{sec:vector} are due to a run of GPT-6 Astra (OpenAI,
ChatGPT Pro mode) that ended on 18 November 2026 (run FREX; in Part~A of its text, the statements on coinvariants, on the
convex Ramsey property, on signed averaging, on the weights for posets, on
ordered posets, and on the amenability, the parity obstruction and the
rational exactness of ordered vector spaces, the last with its two lemmas),
in answer to a research prompt of this project that followed up the
question treated in \cite{D1}. The run built on earlier runs of the same
model, from one of which the first half of Lemma~\ref{lem:ramsey} comes.
The run cites Kechris and Soki\'c for the non-amenability in
Theorem~\ref{thm:posets} and adds a proof of its own, which is not
reproduced here. Due to the writer, Claude Fable~5.2 (Anthropic) at effort
max, on 25 November 2026: the direct proof of
Proposition~\ref{prop:affine}, for which the run cites Isbell and Mitchell;
the formulation of Theorem~\ref{thm:transfer} by the sets $E(B)$; the
second sentence of Theorem~\ref{thm:dict} and Remark~\ref{rem:topos}, after
a remark in the project's notes by Claude Opus~5.5 at effort max;
Theorem~\ref{thm:vector}\,(c) over every ring in which $2$ is not
invertible, by reduction to $R/2R$;
Lemma~\ref{lem:open}, Corollary~\ref{cor:evans} and
Remarks~\ref{rem:zucker}, \ref{rem:HS} and~\ref{rem:reducts}.

\paragraph{The text.} Written by Claude Fable~5.2 at effort max on 25
November 2026, in the setting of \cite{D1}. Every proof was derived again
by the writer from the run's text; the proof of Lemma~\ref{lem:ramsey} was
reconstructed from a one-sentence description in the run.

\paragraph{Checks.} The run's text was followed step by step at its intake
by Claude Opus~5.5 at effort max (19 and 24 November 2026). The writer derived
every proof of this paper again and ran a program of its own, in exact
integer arithmetic: it confirms Lemma~\ref{lem:weights} for all $4474$
partial orders on at most five labelled points ($102085$ one-point
identities, the number the run reports), the three-point example of
Remark~\ref{rem:HS}, and the two identities of the proof of
Theorem~\ref{thm:vector}\,(c) for all injective linear maps into $\F_2^n$,
$n\le5$, with the order $\sqsubset$ of the standard basis. The text was checked on 25 November 2026 by an independent agent, Claude Fable~5.2 at effort max: every proof followed; Theorems~\ref{thm:dict} and~\ref{thm:transfer}, Proposition~\ref{prop:affine}, Theorem~\ref{thm:vector}\,(c) and the non-amenability in Theorem~\ref{thm:posets} proved again independently, the last two by other routes; every citation but \cite{Ho} opened; both programs run. Its repairs to Remark~\ref{rem:reducts}, to the proof of Theorem~\ref{thm:transfer}\,(c) and to Remark~\ref{rem:evans} are included. The sentence of Section~\ref{sec:intro} on the translations of $\Z$ was added by the writer after that check. Nothing is formalized.

\paragraph{The search of the literature.} On 25 November 2026 the writer
read \cite{HS} (the definitions, Sections~3, 4, 10, 11, 16 in part),
\cite{Sn} (Sections~1--3, 8 and the statements), \cite{Ev} (Section~1,
Remarks~2.9--2.12, Lemma~3.7), \cite{Ca} (Section~2), \cite{DLLX} (the
passage of Theorem~5.1), \cite{KS} (Remark~3.2) and \cite{Zu} (Section~3),
and searched the titles and abstracts of arXiv, the web and the Hexagon
catalogue for statements relating exactness of colimits or of coinvariants
to amenability; none was found. The search does not cover full texts.

\paragraph{Citations.} Opened by the writer at the places cited: \cite{Ca}, \cite{D1}, \cite{DLLX}, \cite{Ev}, \cite{HS}, \cite{KS}, \cite{Sn}, \cite{Zu}. Opened by the run and by the independent check of 25 November 2026, at the places cited: \cite{Bo} (arXiv and journal versions), \cite{HN}, \cite{IM}, \cite{KPT}. Opened by the independent check: \cite{vN} (p.~305). Not opened: \cite{Ho}; the theorem used is \cite[Theorem~2.2]{KPT}. Artin's theorem is quoted from \cite[Theorem~5.1]{DLLX}, which takes it from Mac Lane and Moerdijk's book; the book was not opened.

\paragraph{Records.} This paper is accompanied in its workbench by the file
\texttt{CHECK\_EVENTS.tsv}, which holds every check with model, effort and
date, the file \texttt{PROVENANCE.md}, which lists every contributor and
the state of every citation, and the program \texttt{frex\_checks.py} in
the folder \texttt{checks}. The workbenches are reached from
\url{https://github.com/KokunoYumeto/mathematics-workbenches}.

\begin{thebibliography}{99}
\raggedright
\bibitem{Bo} M.~Bodirsky, New Ramsey classes from old, Electron.\ J.\
 Combin.\ 21 (2014), no.~2, Paper~2.22.
\bibitem{Ca} O.~Caramello, Fra\"iss\'e's construction from a topos-theoretic
 perspective, arXiv:0805.2778v3; Log.\ Univers.\ 8 (2014), 261--281.
\bibitem{D1} Exact, total and filtered small categories, Hexagon
 2611.01026, 25 November 2026 (AI-produced: runs of GPT-6 Astra, ChatGPT
 Pro mode; text by Claude Fable~5.2 at effort max).
\bibitem{DLLX} Z.~Di, L.~Li, L.~Liang and F.~Xu, Sheaves of modules on
 atomic sites and discrete representations of topological groups,
 arXiv:2108.13600v4.
\bibitem{Ev} D.~M.~Evans, Permutation modules for Ramsey structures,
 arXiv:2603.29606v2, 9 September 2026.
\bibitem{HS} N.~Harman and A.~Snowden, Oligomorphic groups and tensor
 categories, arXiv:2204.04526v2; Invent.\ Math.\ (2026),
 doi:10.1007/s00222-026-01452-2.
\bibitem{Ho} W.~Hodges, \emph{Model theory}, Cambridge University Press,
 1993.
\bibitem{HN} J.~Hubi\v cka and J.~Ne\v set\v ril, All those Ramsey classes
 (Ramsey classes with closures and forbidden homomorphisms), Adv.\ Math.\
 356 (2019), 106791; arXiv:1606.07979v4.
\bibitem{IM} J.~Isbell and B.~Mitchell, Exact colimits and fixed points,
 Trans.\ Amer.\ Math.\ Soc.\ 220 (1976), 289--298.
\bibitem{KPT} A.~S.~Kechris, V.~G.~Pestov and S.~Todorcevic, Fra\"iss\'e
 limits, Ramsey theory, and topological dynamics of automorphism groups,
 Geom.\ Funct.\ Anal.\ 15 (2005), 106--189.
\bibitem{KS} A.~S.~Kechris and M.~Soki\'c, Dynamical properties of the
 automorphism groups of the random poset and random distributive lattice,
 Fund.\ Math.\ 218 (2012), 69--94.
\bibitem{Sn} A.~Snowden, Measures on partial orders, arXiv:2609.02021v1,
 2 September 2026.
\bibitem{vN} J.~von Neumann, Zur Theorie der Gesellschaftsspiele, Math.\
 Ann.\ 100 (1928), 295--320.
\bibitem{Zu} A.~Zucker, Amenability and unique ergodicity of automorphism
 groups of Fra\"iss\'e structures, Fund.\ Math.\ 226 (2014), 41--61.
\end{thebibliography}
\end{document}
